% Termo de Autorização para Publicação Eletrônica no RI da UFAPE
% X marca a caixa; {} deixa em branco.
\documentclass[10pt,a4paper]{article}

% BLOCO DE DADOS: ALTERE APENAS OS CAMPOS ABAIXO
% Material bibliográfico, X marca a caixa
% Monografia
\newcommand{\FMonografia}{}
% Artigo científico
\newcommand{\FArtigo}{}
% ESO
\newcommand{\FESO}{}
% Dissertação
\newcommand{\FDissertacao}{}

% Autor e documento
% Nome completo, ex.: Maria da Silva
\newcommand{\FAutor}{}
% CPF, ex.: 000.000.000-00
\newcommand{\FCPF}{}
% DDD, vai no parênteses, ex.: 87
\newcommand{\FDDD}{}
% Telefone sem o DDD, ex.: 9 9999-9999
\newcommand{\FFone}{}
% E-mails, ex.: nome@exemplo.com
\newcommand{\FEmails}{}
% Curso, ex.: Ciência da Computação
\newcommand{\FCurso}{}
% Título inteiro, ex.: Estudo sobre revisões sistemáticas
\newcommand{\FTitulo}{}
% Orientador, ex.: Prof. Nome Exemplo
\newcommand{\FOrientador}{}
% Número de folhas, ex.: 120
\newcommand{\FFolhas}{}
% Anexo, X marca sim
\newcommand{\FAnexoSim}{}
% Anexo, X marca não
\newcommand{\FAnexoNao}{}
% Data de defesa, ex.: 10 de Janeiro de 2026
\newcommand{\FData}{}

% Acesso ao documento
% Liberação total, X marca
\newcommand{\FLibTotal}{}
% Liberação parcial, X marca
\newcommand{\FLibParcial}{}
% Permissão sumário, X marca
\newcommand{\FPermSumario}{}
% Permissão capítulos, X marca
\newcommand{\FPermCapitulos}{}
% Quais capítulos, ex.: 2 e 3
\newcommand{\FQuaisA}{}
% Permissão outros itens, X marca
\newcommand{\FPermOutros}{}
% Quais outros, ex.: apêndices
\newcommand{\FQuaisB}{}
% Justificar as restrições, ex.: patente em andamento
\newcommand{\FJust}{}
% Patente, X marca sim
\newcommand{\FPatSim}{}
% Patente, X marca não
\newcommand{\FPatNao}{}
% Se sim, qual patente, ex.: BR 0000000-0
\newcommand{\FPatQual}{}

% Local e data, vão na última linha
% Cidade, ex.: Garanhuns
\newcommand{\FCidade}{Garanhuns}
% Dia, em número, ex.: 10
\newcommand{\FDia}{}
% Mês, por extenso, ex.: Janeiro
\newcommand{\FMes}{}
% Ano, em número, ex.: 2026
\newcommand{\FAno}{}

\usepackage[T1]{fontenc}
\usepackage{mathptmx}
\usepackage[brazilian]{babel}
\usepackage{graphicx}
\usepackage{xcolor}
\usepackage{etoolbox}
\usepackage{amssymb}
\usepackage{amsmath}
\usepackage{ragged2e}
\usepackage[normalem]{ulem}
\renewcommand{\ULthickness}{0.4pt}
\usepackage[a4paper,left=1.27cm,right=1.34cm,top=1.60cm,bottom=0.30cm]{geometry}
\usepackage[bookmarksnumbered=true,psdextra,pdfencoding=auto]{hyperref}
\hypersetup{hidelinks}

\pagestyle{empty}
\setlength{\parindent}{0pt}
\setlength{\parskip}{0pt}

\newcommand{\gap}[2]{\vspace{#2mm}}

\newcommand{\Uaux}[2]{\uline{#1\hspace*{#2}}}
\newcommand{\U}[2]{\expandafter\Uaux\expandafter{#2}{#1}}

\newcommand{\UFaux}[1]{\uline{#1\hspace*{\fill}}}
\newcommand{\UF}[1]{\expandafter\UFaux\expandafter{#1}}

\newcommand{\Uauto}[1]{%
  \ifdefempty{#1}{\U{1.8cm}{}}{\uline{#1}}}

\newcommand{\caixa}[1]{%
  \raisebox{0.9pt}{$\vcenter{\hbox{%
    \setlength{\fboxsep}{1pt}\setlength{\fboxrule}{0.6pt}%
    \ifdefempty{#1}%
      {\fcolorbox{black}{white}{\phantom{{\large$\pmb{\checkmark}$}}}}%
      {\fcolorbox{black}{black!55}{\color{white}{\large$\pmb{\checkmark}$}}}%
  }}$}}

\newcommand{\flinha}[2]{\noindent#1\hspace{0.15cm}\UF{#2}\par}

\newsavebox{\boxC}
\newcommand{\dupla}[4]{%
  \sbox{\boxC}{#1\hspace{0.15cm}#2}%
  \ifdim\wd\boxC>0.48\textwidth
    \noindent#1\hspace{0.15cm}\UF{#2}\par\gap{dupla}{1.3}%
    \noindent#3\hspace{0.15cm}\UF{#4}\par
  \else
    \noindent\parbox[t]{0.485\textwidth}{#1\hspace{0.15cm}\UF{#2}}\hfill\parbox[t]{0.485\textwidth}{#3\hspace{0.15cm}\UF{#4}}\par
  \fi}

\newcommand{\obsitem}[1]{\parshape=2 0pt 18.39cm 0.39cm 18.00cm\noindent#1\par}

\begin{document}

{\setlength{\unitlength}{1mm}%
\begin{picture}(0,0)
  \put(0,-10.6){\includegraphics[height=2.20cm]{figuras/brasao-ufape.pdf}}%
  \put(161.3,-10.6){\includegraphics[height=2.20cm]{figuras/sib-ufape.pdf}}%
\end{picture}}%
\fontsize{10}{10.9}\selectfont
\noindent\makebox[\textwidth][c]{%
  \begin{minipage}{14.2cm}\centering\bfseries\fontsize{11}{18.1}\selectfont
    Ministério da Educação\\
    Universidade Federal do Agreste de Pernambuco\\
    Sistema Integrado de Bibliotecas
  \end{minipage}%
}

\gap{cab-titulo}{3.5}
\noindent\makebox[\textwidth][c]{%
  \begin{minipage}{\textwidth}\centering\bfseries\fontsize{12}{15}\selectfont
    Termo de Autorização para Publicação Eletrônica no Repositório Institucional da UFAPE
  \end{minipage}%
}

\gap{titulo-1}{4.5}

\pdfbookmark[1]{1 - Identificacao do Material Bibliografico}{sec1}
\noindent\textbf{1- Identificação do Material Bibliográfico:}
\caixa{\FMonografia}\ Monografia \caixa{\FArtigo}\ Artigo Científico \caixa{\FESO}\ ESO \caixa{\FDissertacao}\ Dissertação

\gap{1-2}{2.5}

\pdfbookmark[1]{2 - Identificacao do Autor(a) e Documento}{sec2}
\noindent\textbf{2- Identificação do Autor(a) e Documento:}

\gap{2-autor}{1.3}
\dupla{Autor(a):}{\FAutor}{CPF:}{\FCPF}

\gap{autor-fone}{1.3}
\dupla{Telefones: (\FDDD)}{\FFone}{E-mails:}{\FEmails}

\gap{fone-curso}{1.3}
\dupla{Curso:}{\FCurso}{Orientador(a):}{\FOrientador}

\gap{curso-titulo}{1.3}
\flinha{Título do Trabalho:}{\FTitulo}

\gap{titulo-num}{1.3}
\noindent Número de folhas:\hspace{0.15cm}\Uauto{\FFolhas}\hspace{0.46cm}Possui Anexo? \caixa{\FAnexoSim}\ sim \caixa{\FAnexoNao}\ não\hspace{0.26cm}Data de Defesa:\hspace{0.15cm}\UF{\FData}\par

\gap{num-3}{3.0}

\pdfbookmark[1]{3 - Informacoes de Acesso ao Documento}{sec3}
\noindent\textbf{3- Informações de Acesso ao Documento:}

\gap{3-lib}{1.5}
\noindent Liberação para publicação:\hspace{0.88cm}%
\textbf{\caixa{\FLibTotal}\ Total\hspace{0.18cm}\caixa{\FLibParcial}\ Parcial\ - }%
\textcolor{red}{\textbf{Observações: A, B e C}}

\gap{lib-em1}{1.0}

\noindent
  Em caso de publicação parcial, assinalar as permissões: \caixa{\FPermSumario}\ sumário;\par

\gap{em1-em2}{0.8}
\flinha{\caixa{\FPermCapitulos}\ capítulos, informe quais:}{\FQuaisA}

\gap{em2-em3}{0.8}
\flinha{\caixa{\FPermOutros}\ outros itens, informe quais:}{\FQuaisB}

\gap{em3-just}{0.8}
\flinha{Em caso de publicação parcial, justificar as restrições:}{\FJust}

\gap{just-pat1}{0.8}
\noindent Este trabalho ocasionará registro de patente? \caixa{\FPatSim}\ sim \caixa{\FPatNao}\ não.\par

\gap{pat1-pat2}{0.5}
\flinha{Se sim, qual?}{\FPatQual}

\gap{pat2-obs}{0.6}

\pdfbookmark[1]{Observacoes A, B e C}{secObs}
\noindent\textbf{Observações}
\gap{obs-a}{-1.6}
\begin{justify}
\obsitem{a) É imprescindível o envio do arquivo em formato digital do trabalho completo, mesmo em se tratando de disponibilização parcial. Os formatos aceitos são: texto (PDF); som (WAV, MPEG, AIFF ou SND), vídeo (MPEG, AVI ou QT).}
\obsitem{b) A restrição da publicação parcial poderá ser mantida por até 12 meses a partir da data de autorização deste termo. O resumo e os metadados ficarão sempre disponibilizados, mas o arquivo do trabalho apenas será disponibilizado após o fim da restrição. A restrição pode ser prorrogada por mais 12 meses com uma nova solicitação à Biblioteca;}
\obsitem{c) No caso de trabalho com patente, a restrição será total. Significa que o resumo, metadados e arquivo do trabalho serão publicados \textbf{apenas decorrido o período de 24 meses} a partir da data da assinatura deste termo, ou após recebimento de declaração do orientador de que a patente foi registrada.}
\end{justify}

\gap{obs-4}{0.8}

\pdfbookmark[1]{4 - Licenca de Permissao de Uso}{sec4}
\noindent\textbf{4- Licença de Permissão de Uso:}
\gap{4-p1}{-1.6}
\begin{justify}
De acordo com a Portaria n.\ 13 de 15 de fevereiro de 2006, da CAPES, e a Lei n.\ 9610/98, e na qualidade de titular dos direitos de autor desta publicação, autorizo a Universidade Federal do Agreste de Pernambuco disponibilizar gratuitamente, esta obra no Repositório Institucional da UFAPE ou em outras plataformas digitais do Sistema Integrado de Bibliotecas da UFAPE, conforme a licença Creative Commons Atribuição 4.0 (\textquotedbl Licença Pública\textquotedbl) - (CC BY) e também ao IBICT, por meio do site bdtd.ibict.br. Não acarretando ressarcimento dos direitos autorais, o texto integral da obra supracitada, conforme permissões assinaladas, para fins de leitura, impressão e/ou \textit{download}, a título de divulgação da produção científica brasileira e declaro para os devidos fins que o trabalho foi corrigido de acordo com as recomendações da banca examinadora e o(a) orientador(a), sendo assim a versão final pode ser entregue e publicada a partir desta data.
\end{justify}

\gap{4-5}{0.8}

\pdfbookmark[1]{5 - Declaracao de Distribuicao nao-exclusiva}{sec5}
\noindent\textbf{5- Declaração de Distribuição não-exclusiva:}
\gap{5-ref}{1.0}
\begin{justify}
O referido autor:\\
a) Declara que o documento entregue é seu trabalho original, e que detém o direito de conceder os direitos contidos nesta licença. Declara também que a entrega do documento não infringe, tanto quanto lhe é possível saber, os direitos de qualquer outra pessoa ou entidade, não acarretando plágio de qualquer outra obra.\\
b) Se o documento entregue contém material do qual não detém os direitos de autor, declara que obteve autorização do detentor dos direitos de autor para conceder à Universidade Federal do Agreste de Pernambuco os direitos requeridos por esta licença, e que esse material cujos direitos são de terceiros está claramente identificado e reconhecido no texto ou conteúdo do documento entregue.\\
c) Se o documento entregue é baseado em trabalho financiado ou apoiado por outra instituição que não a Universidade Federal do Agreste de Pernambuco, declara que cumpriu quaisquer obrigações exigidas pelo respectivo contrato ou acordo.
\end{justify}

\gap{decl-ass}{12.5}

\pdfbookmark[1]{Assinaturas}{secAss}
\noindent
\begin{minipage}[t]{6.67cm}
  \centering
  \rule{6.67cm}{0.4pt}\\
  \textbf{\small Assinatura do Autor(a)}
\end{minipage}%
\hfill
\begin{minipage}[t]{7.22cm}
  \centering
  \rule{7.22cm}{0.4pt}\\
  \textbf{\small Assinatura do Orientador(a) (para anuência)}
\end{minipage}

\gap{ass-local}{10.0}
\noindent\makebox[\textwidth][c]{%
  \FCidade, \FDia\ de \FMes\ de \FAno.}

\end{document}
